% Gaps (repris de gaps.md, 2e passe du 28/09/2026) et questions de thèse
\section{Zones non couvertes et pistes de recherche}
\label{sec:gaps}

Les gaps ci-dessous viennent de \texttt{gaps.md} (2\textsuperscript{e} passe, 28/09/2026,
73 fiches et 324 lignes). Un gap \emph{étayé} est constaté dans plusieurs fiches (mesure
absente, limite reconnue par les auteurs) ; cela ne prouve pas qu'il est absent de la
littérature. Un gap à l'état d'\emph{hypothèse} ne repose que sur l'absence de papiers
dans la base.

\begin{table}[!htbp]
\centering\footnotesize
\caption{Gaps du corpus drone : statut et piste principale.}
\label{tab:gaps}
\begin{tabularx}{\textwidth}{@{}P{4.3cm}P{1.8cm}L@{}}
\toprule
Gap & Statut & Piste \\
\midrule
Décalage algorithme / matériel : \gl{YOLO}{YOLO} drone (P2, \gl{attention}{attention}) jamais portés sur \gl{FPGA}{FPGA} & étayé &
  YOLOv8n-P2 ou Drone-YOLO nano en \gl{INT8}{INT8} sur KV260 (Vitis~AI), contre \gl{FINN}{FINN} ou hls4ml ; \gl{mAP}{mAP}
  VisDrone avant et après \\
Quantification et petits objets : 0 ligne avec \gl{APS}{AP$_S$} sur 324 & étayé & balayage \gl{PTQ}{PTQ} / \gl{QAT}{QAT}
  (\gl{FP32}{FP32}, INT8, INT4) d'un YOLO drone sur VisDrone, \gl{AP}{AP}$_S$ mesurée sur la carte \\
Tête P2 et haute résolution : coût mémoire on-chip non mesuré & étayé (coût) ; hypothèse
  (absence FPGA) & tailles d'activation par couche de YOLOv8n et YOLOv8n-P2 face au budget
  \gl{BRAM}{BRAM} d'un ZU3EG ou d'une KV260 \\
Paramètres et \gl{FLOPs}{FLOPs} $\neq$ latence sur cible & étayé & trois variantes légères d'un même
  YOLO profilées sur KV260 et Orin, à précision égale \\
Post-traitement sur le \gl{PS}{PS} : goulot de \gl{bout}{bout en bout} & étayé & \gl{NMS}{NMS} matériel ou tête sans NMS ;
  latence caméra $\rightarrow$ sortie instrumentée \\
Calcul adaptatif ou épars jamais évalué à bord & étayé & SAHI à tuiles fixes contre saut de
  tuiles vides sur KV260, latence au pire cas \\
Métriques non comparables ; benchmarks sans métrique d'efficacité & étayé & protocole
  commun : métrique, résolution, \gl{portee}{portée} de la puissance \\
Métrique système absente des papiers YOLO/FPGA & étayé & designs FPGA publiés placés sur le
  \gl{roofline}{modèle F-1} avec masse de carte mesurée \\
Fiabilité des accélérateurs de détection sur drone & étayé (corpus mince) & injection de
  fautes sur un détecteur FPGA avec métrique de mission \\
Pas de comparaison entre \gl{gamme}{gammes de puissance} & étayé & même réseau, même entrée, même
  portée de puissance sur GAP8, FPGA et Jetson \\
Détection et navigation sur le même budget & étayé & partage d'un \gl{SoC}{SoC} par reconfiguration
  partielle, interférence mesurée \\
Suivi embarqué sans puissance ni FPGA & étayé & HiFT ou SiamAPN++ sur FPGA, énergie par
  image \\
Contre-drone évalué sur \gl{GPU}{GPU} de bureau & étayé & YOLOv8n sur Det-Fly, INT8 sur KV260 contre
  Orin, puissance à la carte \\
Détecteur \gl{airair}{air $\rightarrow$ air} embarqué sur intercepteur & étayé & détecteur sur SoC avec
  AP, \gl{FPS}{FPS}, W et latence de bout en bout ; Det-Fly et ARD-MAV \\
Air $\rightarrow$ air : coût du mouvement multi-images & étayé (coût) ; hypothèse (FPGA) &
  profil par étage de GL-YOMO ou GLAD ; tampon on-chip de 8 images \\
Fusion visible + thermique à bord & hypothèse & recherche ciblée (\texttt{/paper-discover}) \\
Apprentissage continu à bord du drone & hypothèse & base \texttt{fpga} ; recherche ciblée \\
Code et bitstreams rarement publiés & étayé (relevé approximatif) & distinguer code
  d'entraînement et design matériel \\
\bottomrule
\end{tabularx}
\end{table}

\subsection{Trois gaps structurants}

\paragraph{Le décalage algorithme / matériel.} C'est l'hypothèse de départ de la base, et
le corpus l'étaye. Les détecteurs qui font l'état de l'art sur VisDrone ajoutent une tête
P2 et de l'attention (section~\ref{sec:detection}). Les implantations FPGA restent sur des
YOLO anciens et, pour tenir, retirent la branche multi-échelle
\citep{2025-cali-optimized-implementation-yolov3-tiny-real} ou remplacent l'attention
\citep{arxiv-2607.13106-no-attention-no-problem-dpu}\pp. La seule mesure embarquée d'une
\gl{teteP2}{tête P2} (Jetson TX2) montre un facteur 3 à 4 en latence
\citep[Table~6]{2023-zhang-drone-yolo-efficient-neural-network}. Deux familles
d'alternatives à P2 existent, jamais évaluées sur FPGA : transférer l'information de P2
vers P3 (TPH-YOLOv5++), ou recadrer à pleine résolution autour des régions utiles (SAHI,
UFPMP-Det, GLAD).

\paragraph{La précision sur petits objets après quantification.} Aucune des 324 lignes ne
rapporte d'AP$_S$. Les pertes de \gl{quantif}{quantification} connues sont globales et souvent fortes
(6 à 30 points de mAP selon le modèle et le jeu). On ne sait donc pas si la
quantification, indispensable sur FPGA et \gl{MCU}{MCU}, détruit précisément ce que les YOLO drone
cherchent à gagner.

\paragraph{La métrique de mission.} Le lien entre latence, masse et \gl{vitesse}{vitesse sûre} est
modélisé \citep{arxiv-2204.10898-roofline-model-uavs-bottleneck-analysis}\pp{} et mesuré
pour des nano-drones, jamais pour un détecteur sur FPGA. Le modèle lui-même ignore la
qualité de perception, et la gamme nano montre qu'un réseau plus rapide peut moins bien
servir la mission \citep[Table~IV]{arxiv-2103.10873-fully-onboard-ai-powered-human}\pp.

\subsection{Questions de thèse candidates}

Les trois questions ci-dessous sont des déductions à partir des gaps, pas des
conclusions du corpus.
\begin{enumerate}
  \item \textbf{Détecteur air $\rightarrow$ air sur SoC embarquable.} Peut-on détecter un
    drone de quelques pixels à bord d'un intercepteur, sur KV260 ou Zynq UltraScale+, avec
    une AP, une latence de bout en bout et une puissance à la carte mesurées, et comment
    se compare-t-on à une Jetson Orin à protocole égal ? Données : Det-Fly, ARD-MAV,
    NPS-Drones. Point de départ : l'architecture globale/locale de GLAD, dont la partie
    non neuronale est à accélérer.
  \item \textbf{Petits objets et précision numérique.} Quelle est la perte d'AP$_S$ d'un
    YOLO drone (avec et sans tête P2) quand on passe de FP32 à INT8 puis INT4, en PTQ et en
    QAT, mesurée sur la carte ? Quelle alternative à P2 (recadrage, transfert vers P3) tient
    dans la mémoire on-chip d'une petite FPGA à précision égale ?
  \item \textbf{Du FPS à la vitesse sûre.} Comment traduire le débit, la puissance et la
    masse d'un accélérateur de détection en vitesse de vol sûre et en autonomie, en
    intégrant la qualité de détection au modèle F-1 ? À quelle gamme de puissance
    (nano, FPGA, Jetson) correspond l'optimum pour un drone donné ?
\end{enumerate}
Laboratoires français à approcher : \texttt{labs.md} ne recense encore que le LIRMM (test
et fiabilité), sans papier du corpus à l'appui ; la liste reste à construire.
